\newpage
\section{Đánh giá kết quả}

Sau khi tiến hành thực nghiệm với 1 000 000 lần tăng giá trị trên hệ thống, các số liệu thu được đã phản ánh rất rõ bản chất và giới hạn của từng cơ chế đồng bộ hóa khi đối mặt với áp lực tranh chấp gia tăng. Dưới đây là bảng tổng hợp kết quả thời gian thực thi (giây) theo số lượng luồng:








\begin{table}[h!]
    \centering

    \resizebox{\textwidth}{!}{%
    \begin{tabular}{|l|l|l|l|l|l|l|l|l|l|l|}
    \hline
    Threads & Mutex & Atom(Seq) & Atom(Rel) & CAS(Weak) & CAS(Str) & CAS(Wk+BO) & CAS(Str+BO) & SpinLock & Ticket & Semaphore \\ \hline
    1 & 0.0274 & 0.025 & 0.0159 & 0.0085 & 0.0083 & 0.0087 & 0.0084 & 0.0067 & 0.007 & 0.1432 \\ \hline
    6 & 0.0512 & 0.0278 & 0.0237 & 0.0437 & 0.0438 & 0.009 & 0.009 & 0.0304 & 0.2264 & 0.1199 \\ \hline
    11 & 0.0782 & 0.0294 & 0.0297 & 0.0618 & 0.053 & 0.0104 & 0.01 & 0.0329 & 0.3466 & 0.2112 \\ \hline
    16 & 0.088 & 0.0298 & 0.0291 & 0.0591 & 0.0627 & 0.0102 & 0.0103 & 0.0515 & 0.427 & 0.2542 \\ \hline
    21 & 0.0762 & 0.032 & 0.0296 & 0.0621 & 0.0621 & 0.0102 & 0.01 & 0.0383 & 1.2014 & 0.2357 \\ \hline
    26 & 0.0771 & 0.0309 & 0.0278 & 0.0681 & 0.0676 & 0.012 & 0.0117 & 0.0388 & 3.192 & 0.2519 \\ \hline
    31 & 0.0667 & 0.0303 & 0.0263 & 0.0671 & 0.0673 & 0.0121 & 0.0117 & 0.0399 & 2.1302 & 0.2387 \\ \hline
    36 & 0.0771 & 0.0297 & 0.0263 & 0.0701 & 0.0681 & 0.0127 & 0.0125 & 0.0384 & 2.369 & 0.2277 \\ \hline
    41 & 0.0768 & 0.0266 & 0.0278 & 0.0621 & 0.066 & 0.0127 & 0.0128 & 0.0335 & 4.9944 & 0.2339 \\ \hline
    46 & 0.0722 & 0.0313 & 0.027 & 0.0667 & 0.0676 & 0.0133 & 0.0131 & 0.0318 & 7.0449 & 0.2488 \\ \hline
    51 & 0.0721 & 0.03 & 0.0282 & 0.0716 & 0.0715 & 0.0151 & 0.0167 & 0.0364 & 16.0141 & 0.2438 \\ \hline
    56 & 0.0749 & 0.0285 & 0.0276 & 0.0714 & 0.069 & 0.0152 & 0.0154 & 0.0286 & 10.8053 & 0.2399 \\ \hline
    61 & 0.0743 & 0.0285 & 0.0279 & 0.0664 & 0.0656 & 0.0151 & 0.0153 & 0.0249 & 13.2762 & 0.2168 \\ \hline
    66 & 0.0675 & 0.0282 & 0.0266 & 0.0707 & 0.0712 & 0.0165 & 0.0158 & 0.0227 & 13.2848 & 0.2459 \\ \hline
    71 & 0.0718 & 0.0283 & 0.0291 & 0.0683 & 0.0688 & 0.0159 & 0.0162 & 0.0248 & 14.9874 & 0.2502 \\ \hline
    76 & 0.0806 & 0.0339 & 0.0278 & 0.0712 & 0.0711 & 0.0169 & 0.017 & 0.0549 & 8.8155 & 0.2513 \\ \hline
    81 & 0.0797 & 0.0313 & 0.0274 & 0.0738 & 0.0708 & 0.0178 & 0.0176 & 0.0247 & 14.9077 & 0.2293 \\ \hline
    86 & 0.0671 & 0.0283 & 0.0255 & 0.0659 & 0.0653 & 0.0169 & 0.0196 & 0.0268 & 14.6838 & 0.2375 \\ \hline
    91 & 0.0605 & 0.0273 & 0.0277 & 0.0682 & 0.0655 & 0.0169 & 0.0183 & 0.0236 & 19.2949 & 0.2532 \\ \hline
    96 & 0.0821 & 0.0307 & 0.0303 & 0.069 & 0.073 & 0.0195 & 0.019 & 0.0278 & 14.5954 & 0.2331 \\ \hline
    \end{tabular}%
    }
        \vspace{0.2cm}
        \caption{Kết quả Benchmark thời gian thực thi của các cơ chế (đơn vị: giây)}

\end{table}

\begin{figure}[H]
    \centering
    \includegraphics[width=0.8\textwidth]{001.png}
    \caption{Biểu đồ so sánh hiệu năng của tất cả cơ chế}
    \label{fig:performance_standard}
\end{figure}


\begin{figure}[H]
    \centering
    \includegraphics[width=0.8\textwidth]{002.png}
    \caption{Biểu đồ so sánh hiệu năng của tất cả cơ chế trừ Ticker Lock}
    \label{fig:performance_ticket}
\end{figure}

\begin{figure}[H]
    \centering
    \includegraphics[width=0.8\textwidth]{003.png}
    \caption{Biểu đồ so sánh hiệu năng của tất cả cơ chế trừ Ticker Lock và Semaphore}
    \label{fig:performance_ticket}
\end{figure}

Khi quan sát kỹ vào bảng số liệu và các biểu đồ thực nghiệm, chúng tôi rút ra những nhận xét quan trọng sau:

\begin{itemize}
       \item \textbf{Mutex:} Biểu đồ cho thấy Mutex duy trì độ ổn định rất cao, dao động trong khoảng $0.06s - 0.08s$ bất chấp số lượng luồng tăng lên. Mặc dù không đạt tốc độ nhanh nhất do context switch của hệ điều hành, nhưng đây là phương án an toàn nhất, đảm bảo tính công bằng và tránh lãng phí CPU khi luồng phải chờ đợi lâu.
    
    \item \textbf{SpinLock:} Cơ chế này có hiệu năng tốt hơn Mutex ở mức tải thấp ($0.03s$) nhưng thể hiện sự thiếu ổn định nghiêm trọng khi số luồng tăng (dao động mạnh từ $0.02s$ đến $0.05s$). Sự trồi sụt này là do SpinLock phụ thuộc hoàn toàn vào bộ lập lịch; nếu luồng đang giữ khóa bị ngắt (preempted), các luồng khác sẽ tiêu tốn tài nguyên CPU vô ích để chờ đợi.
    
    \item \textbf{Ticket Lock:} Đây là trường hợp có thời gian đáng chú ý nhất. Từ mức khởi điểm ấn tượng $0.007s$, thời gian thực thi bùng nổ theo hàm mũ, đạt đỉnh điểm hơn $19s$ tại 91 luồng. Cơ chế "xếp hàng" nghiêm ngặt (FIFO) đã tạo ra nút thắt cổ chai, nơi hàng loạt luồng phải chờ đợi tuần tự, làm tê liệt hoàn toàn hệ thống khi tranh chấp lên cao.

    \item \textbf{Semaphore:} Số liệu thực nghiệm cho thấy đây là cơ chế có chi phí quản lý cao nhất. Ngay từ 1 luồng, thời gian thực thi đã là $0.1432s$, chậm hơn đáng kể so với tất cả các phương pháp khác. Tuy nhiên, nó duy trì được sự ổn định (dao động quanh $0.21s - 0.25s$) khi số luồng tăng cao. Điều này chứng tỏ Semaphore quá "nặng nề" để sử dụng cho mục đích bảo vệ vùng găng nhỏ (như tăng biến đếm). Vai trò thực sự của nó phù hợp hơn cho việc điều phối tín hiệu hoặc quản lý tài nguyên gộp thay vì khóa loại trừ.
\end{itemize}

\begin{itemize}
    \item \textbf{Atomic Relaxed \& Sequential Consistency:} Cả hai phương pháp này đều thể hiện một đường thẳng hiệu năng "phẳng lì" tuyệt vời quanh mức $0.027s$. Sự ổn định này đến từ việc các chỉ thị được ánh xạ trực tiếp xuống phần cứng.
\end{itemize}


\begin{itemize}
    \item \textbf{CAS Nguyên bản (Naive CAS):} Các biến thể \texttt{CAS Weak} và \texttt{CAS Strong} thông thường bộc lộ rõ điểm yếu khi đối mặt với tranh chấp. Thời gian thực thi tăng vọt gấp 10 lần (từ $0.008s$ lên $0.076s$) và bị bão hòa tại đó. Nguyên nhân là do hiện tượng bus saturatio khi hàng loạt luồng liên tục gửi yêu cầu cập nhật bộ nhớ thất bại.
    
    \item \textbf{CAS kết hợp Exponential Backoff:} Đây là điểm sáng nhất của thực nghiệm. Với chiến lược "lùi bước theo hàm mũ", thời gian thực thi được duy trì ở mức siêu thấp ($0.008s - 0.019s$) ngay cả khi chạy với 96 luồng.

\end{itemize}




\textbf{Tổng kết chung:} Kết quả thực nghiệm cho thấy cơ chế \textbf{CAS kết hợp Exponential Backoff} là giải pháp tối ưu nhất về hiệu năng, nhờ khả năng tận dụng tốc độ của nhóm \textbf{Atomic-based} đồng thời giảm thiểu tranh chấp bộ nhớ thông qua chiến thuật chờ đợi thông minh. Các phương pháp \textbf{Lock-based} truyền thống vẫn giữ được ưu thế về độ ổn định và an toàn, trong khi \textbf{Ticket Lock} cần được cân nhắc kỹ lưỡng khi triển khai ở quy mô luồng lớn để tránh gây ra hiện tượng nghẽn băng thông hoặc sụt giảm nghiêm trọng hiệu suất hệ thống.
